\scanpage{21}
\mspage{\handwritten{\unsure{15|35}}}
me $R_s$ est un tore maximal dans $G_s$, et la formation de \uline{Norm}, \uline{Centr}, \uline{Norm}/\uline{Centr} est compatible avec toute extension de la base, en particulier le passage aux fibr\unsure{e|è|é}s, on voit que l'on a
\[
w(s)=w'(s)\qquad ;
\]
je dis que de plus, pour $t$ voisin de $s$, on aura les inégalités
\[
w(t)\handwritten{\geq}w'(t)\qquad\text{et}\qquad w'(t)\handwritten{\geq}w'(s)\quad ,
\]
ce qui suffira pour établir a). \deleted{La deuxième} \handwritten{\unsure{C|L}}es deux inégalités \deleted{\unsure{résulte|résultera}} sont contenu\handwritten{es} dans les deux lemmes suivants:

\uline{Lemme} \handwritten{2}.2. Soi\deleted{\unsure{t}}ent $S$ un préschéma, $W$ un $S$-préschéma en groupes qui est étale, \deleted{quasi-fini et} séparé et quasi-fini sur $S$. Pour tout $s\handwritten{\in}S$, soit $f(s)$ la classe de la fibre géométrique de $W$ en $s$, qui est un \deleted{ensemble}\inserted{élément} de l'ensemble ordonné $E$ des classes de groupes finis à isomorphisme près, introduit dans \handwritten{2}.1. Alors la fonction $f:S\handwritten{\to}E$ est semi-continue \deleted{\unsure{sup}}inférieurement. Pour qu'elle soit constante au voisinage de $s\handwritten{\in}S$\handwritten{,} \deleted{(ce qui revient au même, pour} i\deleted{\unsure{o}}l faut et suffit qu'il en soit de même de la fonction $s\handwritten{\rightsquigarrow}\operatorname{card}f(s)$, et pour ceci il faut et suffit que \handwritten{$W$} soit fini sur $S$ au dessus d'un voisinage ouvert $U$ de $s$.

Ce résultat est un raffinement, en termes de groupes, de celui invoqué dans la démonstration de \inserted{\handwritten{Exp $\overline{\underline{\mathrm{XI}}}$}} \handwritten{5}.10. , nous nous bornons à une esquisse de \inserted{\handwritten{la}} démonstration (\inserted{\handwritten{qui est}} d\handwritten{u} type \handwritten{le} plus standart). On est ramené comme d'habitude au cas $S$ affine noethérien. On voit de suite que la fo\deleted{\unsure{u}}nction $f$ est constructible (EGA \deleted{\unsure{IV|II}} $\unsure{0|O}_{\mathrm{III}}$ 9.3.1. et 9.3.2. , et sor\handwritten{it}es de EGA IV 9), et en vertu de EGA $0_{\mathrm{III}}$ 9.3.4. on est ramené \inserted{pour la semi-continuité} à prouver que si $t$ est une générisation de $s$, on a $f(t)\handwritten{\geq}f(\deleted{\unsure{x}}s)$. Cela nous ramène \deleted{au cas} grâce à EGA II 7.1.7. au cas où $S$ est le spectre d'un anneau de valuation discrète, qu'on peut supposer complet à corps résiduel
% archive sheet number: 20
